\subsection{函子性质}

命题 2. 设 X 与 Y 为两个集合。每个映射 \(u: \mathrm{X} \to \mathrm{Y}\) 都可唯一地扩张为李代数同态 \(\mathrm{L}(u): \mathrm{L}(\mathrm{X}) \to \mathrm{L}(\mathrm{Y})\)。对每个映射 \(v: \mathrm{Y} \to \mathrm{Z}\)，有 \(\mathrm{L}(v \circ u) = \mathrm{L}(v) \circ \mathrm{L}(u)\)。

\(\mathrm{L}(u)\) 的存在性与唯一性由第 2 号命题 1 得出。同态 \(\mathrm{L}(v\circ u)\) 与 \(\mathrm{L}(v)\circ \mathrm{L}(u)\) 在 X 上有相同的限制，因而相等（命题 1）。

推论. 若 \(u\) 是单射（相应地，满射、双射），则 \(L(u)\) 亦然。

由于断言对 \(\mathrm{X} = \varnothing\) 是平凡的，我们设 \(\mathrm{X} \neq \varnothing\)。若 \(u\) 是单射，则存在映射 \(v\)（\(\mathrm{Y}\) 到 \(\mathrm{X}\) 的），使 \(v \circ u\) 是 \(\mathrm{X}\) 的恒同映射；由命题 2，\(\mathrm{L}(v) \circ \mathrm{L}(u)\) 是 \(\mathrm{L}(\mathrm{X})\) 的恒同自同构，因而 \(\mathrm{L}(u)\) 是单射。当 \(u\) 是满射时，存在映射 \(w\)（\(\mathrm{Y}\) 到 \(\mathrm{X}\) 的），使 \(u \circ w\) 是 \(\mathrm{Y}\) 的恒同映射；于是 \(\mathrm{L}(u) \circ \mathrm{L}(w)\) 是 \(\mathrm{L}(\mathrm{Y})\) 的恒同映射，这就证明 \(\mathrm{L}(u)\) 是满射。

设 \(\mathbf{X}\) 为集合，\(\mathbf{S}\) 为 \(\mathbf{X}\) 的子集。上述推论表明，\(\mathbf{S}\) 到 \(\mathbf{X}\) 的典范单射可扩张为同构 \(\alpha\)（\(\mathbf{L}(\mathbf{S})\) 到李子代数 \(\mathbf{L}'(\mathbf{S})\) 上的同构，后者是 \(\mathbf{L}(\mathbf{X})\) 中由 \(\mathbf{S}\) 生成的李子代数）；我们将把 \(\mathbf{L}(\mathbf{S})\) 与 \(\mathbf{L}'(\mathbf{S})\) 借助 \(\alpha\) 等同起来。

设 \((\mathrm{S}_{\alpha})_{\alpha \in \mathbf{I}}\) 为 X 的子集的一个右向族，其并为 S。关系 \(\mathrm{S}_{\alpha} \subset \mathrm{S}_{\beta}\) 蕴含 \(\mathrm{L}(\mathrm{S}_{\alpha}) \subset \mathrm{L}(\mathrm{S}_{\beta})\)，因而李子代数族 \(\mathrm{L}(\mathrm{S}_{\alpha})\)（\(\mathrm{L}(\mathrm{X})\) 的）是右向的。因此 \(\mathfrak{g} = \bigcup_{\alpha \in \mathbf{I}} \mathrm{L}(\mathrm{S}_{\alpha})\) 是 \(\mathrm{L}(\mathrm{X})\) 的李子代数；于是 \(\mathrm{S} \subset \mathfrak{g}\)，由此 \(\mathrm{L}(\mathrm{S}) \subset \mathfrak{g}\)，又由于 \(\mathrm{L}(\mathrm{S}_{\alpha}) \subset \mathrm{L}(\mathrm{S})\) 对一切 \(\alpha \in \mathbf{I}\)，故 \(\mathfrak{g} \subset \mathrm{L}(\mathrm{S})\)。于是

\[
\mathrm{L} \left(\bigcup_ {\alpha \in \mathrm{I}} \mathrm{S} _ {\alpha}\right) = \bigcup_ {\alpha \in \mathrm{I}} \mathrm{L} (\mathrm{S} _ {\alpha})\tag{7}
\]

对 X 的子集的每个右向族 \((\mathrm{S}_{\alpha})_{\alpha\in\mathbb{I}}\) 成立。

把上述结果应用于 X 的有限子集的族，我们看到 \(\mathrm{L}(\mathrm{X})\) 的每个元素都具有形式 \(\mathrm{P}(x_{1},\ldots,x_{n})\)，其中 P 是 n 个不定元的李多项式，而 \(x_{1},\ldots,x_{n}\) 是 X 的元素。

命题 3. 设 \(\mathbf{K}'\) 为非零交换环，\(u: \mathbf{K} \to \mathbf{K}'\) 为环同态。对每个集合 \(\mathbf{X}\)，存在且仅存在一个李 \(\mathbf{K}'\) -代数同态

\[
v \colon \mathrm{L} _ {\mathbf {K}} (\mathrm{X}) \otimes \mathrm{K} ^ {\prime} \rightarrow \mathrm{L} _ {\mathbf {K} ^ {\prime}} (\mathrm{X})
\]

使得 \(v(x \otimes 1) = x\) 对 \(x \in \mathbf{X}\) 成立。而且，\(v\) 是同构。

把命题 1 应用于视为李 K-代数的 \(\mathfrak{g} = \mathrm{L}_{\mathbb{K}'}(\mathrm{X})\) 以及映射 \(x\mapsto x\)（\(\mathbf{X}\) 到 \(\mathfrak{g}\) 的），我们得到一个 K-同态 \(\mathrm{L}_{\mathbb{K}}(\mathrm{X})\to \mathrm{L}_{\mathbb{K}'}(\mathrm{X})\)，从而存在 K'-同态 \(v\colon \mathrm{L}_{\mathbb{K}}(\mathrm{X})\otimes \mathrm{K}'\to \mathrm{L}_{\mathbb{K}'}(\mathrm{X})\)。\(v\) 的唯一性以及它是同构这一事实，由有序对 \((\mathrm{L}_{\mathbb{K}}(\mathrm{X})\otimes \mathrm{K}',x\mapsto x\otimes 1)\) 与有序对 \((\mathrm{L}_{\mathbb{K}'}(\mathrm{X}),x\mapsto x)\) 是同一泛映射问题的解这一事实得出。

注. 设 \(\mathfrak{h}'\) 为李 K'-代数，\(\mathfrak{h}\) 为由 \(\mathfrak{h}'\) 通过限制标量环而得到的李 K-代数。若 \(\mathrm{P} \in \mathrm{L}_{\mathrm{K}}(\mathrm{X})\)，我们可以定义 \(\tilde{\mathrm{P}}: \mathfrak{h}^{\mathrm{X}} \to \mathfrak{h}\)（第 4 号）。我们立即看出

\[
\tilde {\mathrm{P}} = (v (\mathrm{P} \otimes 1)) ^ {\sim}.
\]
% semantic-fragments: legacy-input-seam
\relax{}
